\section*{Clustering, density and outliers}
\subsection*{K-means and centroids}
\textbf{Purpose:} Partition observations into $K$ centre-based clusters.
\[
\min_{\{C_k,\mu_k\}}\sum_{k=1}^K\sum_{i\in C_k}\|x_i-\mu_k\|^2,\qquad
\mu_k=\frac1{|C_k|}\sum_{i\in C_k}x_i.
\]
\textbf{Symbols:} $K$ = number of clusters; $C_k$ = set of rows assigned to cluster $k$; $|C_k|$ = its size; $\mu_k$ = centroid/mean vector; $\|\cdot\|^2$ = squared Euclidean distance minimized by K-means.

\textbf{How to read output:} each observation belongs to its nearest centroid; a centroid table is a mean attribute profile for each cluster.

\textbf{Do this:} compute distances to centroids; assign nearest; recompute cluster means; repeat assignments if requested.

\subsection*{Hierarchical clustering and dendrograms}
\begin{tabularx}{\linewidth}{@{}lX@{}}
\toprule
\textbf{Linkage} & \textbf{Cluster-to-cluster distance}\\ \midrule
Single & minimum pair distance; joins through close links.\\
Complete & maximum pair distance; favours compact groups.\\
Average & average of all cross-cluster pair distances.\\
Centroid & distance between cluster means.\\
\bottomrule
\end{tabularx}
\textbf{How to read a dendrogram:} leaves are observations; branch height is merge distance; a horizontal cut gives the connected groups below that height.

\subsection*{Compare clusterings}
\[
\operatorname{Rand}=\frac{f_{11}+f_{00}}{\binom N2},\qquad
J=\frac{f_{11}}{f_{11}+f_{10}+f_{01}}.
\]
\textbf{Symbols:} $\binom N2$ = number of observation pairs; $f_{11}$ = together in both clusterings; $f_{00}$ = apart in both; $f_{10},f_{01}$ = pair-status disagreements; $J$ = Jaccard agreement.

\textbf{Means:} Rand counts both kinds of agreement; Jaccard measures overlap of same-cluster pairs only.

\subsection*{Gaussian mixture model}
\textbf{Purpose:} Describe data by weighted Gaussian components with soft memberships.
\[
\begin{gathered}
p(x)=\sum_{k=1}^K\pi_k\mathcal N(x|\mu_k,\Sigma_k),
\qquad \sum_k\pi_k=1,\\
\gamma_{ik}=\frac{\pi_k\mathcal N(x_i|\mu_k,\Sigma_k)}
{\sum_j\pi_j\mathcal N(x_i|\mu_j,\Sigma_j)} .
\end{gathered}
\]
\textbf{Symbols:} $p(x)$ = mixture density; $\pi_k$ = weight of component $k$; $\mathcal N(x|\mu_k,\Sigma_k)$ = its Gaussian density; $\gamma_{ik}$ = responsibility for row $i$; $\mu_k,\Sigma_k$ = component mean and covariance; $j$ in the denominator indexes components.
\textbf{How to read output:} $\pi_k$ is component frequency; $\gamma_{ik}$ is observation $i$'s responsibility for component $k$ and sums to $1$ across components. Largest responsibility gives a hard label.

\subsection*{KDE and local density}
\textbf{Purpose:} Find regions of low probability density and potential outliers.
\[
\hat p(x)=\frac1N\sum_i\mathcal N(x|x_i,\sigma^2I),
\]
\[
\ell_{\rm LOO}(\sigma)=\sum_i\log\!\left[
\frac1{N-1}\sum_{j\ne i}\mathcal N(x_i|x_j,\sigma^2I)\right].
\]
\[
\operatorname{dens}(x_i,K)=
\left(\frac1K\sum_{x_j\in N_i}d(x_i,x_j)\right)^{-1},
\]
\[
\operatorname{ARD}(x_i,K)=
\frac1K\sum_{x_j\in N_i}
\frac{\operatorname{dens}(x_i,K)}{\operatorname{dens}(x_j,K)}.
\]
\textbf{Symbols:} $\hat p$ = kernel density estimate; $\sigma$ = Gaussian bandwidth; $I$ = identity matrix; $\ell_{\rm LOO}$ = leave-one-out log likelihood; $N_i$ = $K$ nearest neighbours of $x_i$; $\operatorname{dens}$ = inverse mean neighbour distance; ARD = average relative density ratio.
\textbf{How to read it:} small KDE/local density marks isolation; small ARD marks a point less dense than its neighbours. Larger $\sigma$ smooths a KDE curve.

\textbf{Do this:} exclude the evaluated point for leave-one-out/local-neighbour calculations; compute densities; identify the smallest density or ARD as most outlying.
